% ---------------------------------------------------------------------------
% Preambule du livrable ST3.2 -- ARD JUNON, projet Prediction
% Compilation : pdflatex + biber (voir Makefile a la racine, cible st32)
%
% Base sur le gabarit fourni par N. Labroche. Deux ecarts assumes :
%   - biblatex/biber au lieu de bibtex, pour partager bibliographie.bib avec
%     le rapport de cloture sans dupliquer d'entrees ;
%   - quelques paquets ajoutes (siunitx, booktabs, tabularx, float, enumitem)
%     requis par la section des developpements techniques.
% ---------------------------------------------------------------------------

% ---------------------------------------------------------------------------
% CHAMPS ADMINISTRATIFS
%
% Seule la STRUCTURE de la page de garde vient du gabarit consulte (un livrable
% d'une autre equipe) : les libelles de formulaire, l'ordre des blocs, le
% tableau des auteurs, les deux logos. Aucun contenu de ce document n'est repris.
%
% ATTENTION, une valeur fait exception et doit etre verifiee avant remise :
% le numero de convention ci-dessous a ete RECOPIE de ce gabarit. Il identifie
% le projet PREDICTION, donc il s'applique vraisemblablement aussi a nous, mais
% cela n'a pas ete verifie sur le texte de la convention. A confirmer.
%
% Deux autres champs sont deduits de ce que le present document declare de
% lui-meme, et non d'une source administrative :
%   - le numero de livrable (le gabarit numerote « L1.1 », ce document s'annonce
%     « ST3.2 » : verifier lequel la convention emploie) ;
%   - l'organisme porteur du livrable.
%
% La date doit etre figee a la remise, et non laissee a \today. Et les prenoms
% de deux coauteurs manquent : voir sections/00-couverture.tex.
% ---------------------------------------------------------------------------
\newcommand{\projetLivrable}{PREDICTION -- ARD CVL JUNON}
\newcommand{\conventionLivrable}{Convention d'Application
  N\textsuperscript{o}2021 00146736 / N\textsuperscript{o}2022 00151803}
\newcommand{\intituleLivrable}{ST3.2 Explicabilité et interprétabilité dans les
  modèles prédictifs profonds de séries temporelles}
\newcommand{\organismeLivrable}{Université de Tours, LIFAT}
\newcommand{\dateLivrable}{\today}
\newcommand{\titreCourant}{Explicabilité et interprétabilité dans les modèles prédictifs
  profonds de séries temporelles}

% --- PAQUETS DE BASE ---
\usepackage[utf8]{inputenc}
\usepackage[french]{babel}
\usepackage[T1]{fontenc}
\usepackage{lmodern}
\usepackage{microtype}

% --- GEOMETRIE ET MISE EN PAGE ---
\usepackage[top=2.5cm, bottom=2.5cm, left=2.5cm, right=2.5cm]{geometry}
\usepackage{setspace}
\onehalfspacing % Interligne 1.5

% --- GRAPHISME ET COULEURS ---
\usepackage{graphicx}
\usepackage{xcolor}
\definecolor{primary}{RGB}{24, 76, 120}   % Bleu académique
\definecolor{secondary}{RGB}{100, 110, 120} % Gris neutre
\definecolor{accent}{RGB}{180, 40, 40}    % Accentuation

% --- TABLEAUX ET LISTES ---
\usepackage{booktabs}
\usepackage{tabularx}
\usepackage{array}
\usepackage{float}
\renewcommand{\arraystretch}{1.15}
\newcolumntype{L}[1]{>{\raggedright\arraybackslash}p{#1}}
\newcolumntype{Y}{>{\raggedright\arraybackslash}X}

\usepackage{enumitem}
\setlist[itemize]{leftmargin=1.4em,itemsep=0.25em,topsep=0.4em}
\setlist[enumerate]{leftmargin=1.6em,itemsep=0.25em,topsep=0.4em}

% --- FIGURES ---
\usepackage{tikz}
\usetikzlibrary{arrows.meta,positioning,shapes.geometric,fit,backgrounds,calc}

% --- LEGENDES ---
\usepackage{caption}
\captionsetup{font=small,labelfont={bf,color=primary}}

% --- UNITES ---
\usepackage{siunitx}
\sisetup{
  output-decimal-marker = {,},
  group-separator       = {\,},
  group-minimum-digits  = 4,
  range-phrase          = {\ à\ },
  range-units           = single,
}
% L'annee n'est pas une unite siunitx : on la declare.
\DeclareSIUnit{\an}{an}

% --- BIBLIOGRAPHIE ---
% Deux ressources : le fonds commun aux deux documents du depot, et les
% references propres au ST3.2. Aucune entree n'est dupliquee entre les deux.
\usepackage[backend=biber,style=numeric,sorting=nyt,maxbibnames=6,giveninits=true]{biblatex}
\addbibresource{../bibliographie.bib}
\addbibresource{references.bib}

% --- HYPERLIENS ---
\usepackage{hyperref}
\hypersetup{
    colorlinks=true,
    linkcolor=primary,
    citecolor=primary,
    urlcolor=primary,
    pdftitle={ST3.2 Prédiction - ARD JUNON},
    pdfauthor={N. Labroche, N. Ringuet, E. Doumard, M. Gol Pour},
    pdflang={fr-FR},
    pdfsubject={Explicabilité et interprétabilité dans les modèles prédictifs profonds de séries temporelles},
    pdfkeywords={XAI, jumeau numérique, hydrogéologie, séries temporelles multivariées, explications contrefactuelles, attribution de caractéristiques},
}

% Les URL longues de la bibliographie ne se coupent pas d'elles-memes et
% debordent de la justification : xurl autorise la cesure a tout caractere.
\usepackage{xurl}

% --- EN-TÊTES ET PIEDS DE PAGE ---
% Le gabarit porte le titre du livrable en en-tete et « Page N sur M » en pied.
\usepackage{lastpage}
\usepackage{fancyhdr}
\pagestyle{fancy}
\fancyhf{}
\lhead{\small \color{secondary} \titreCourant}
\rhead{\small \color{secondary} ST3.2}
\cfoot{\small \color{secondary} Page \thepage{} sur \pageref{LastPage}}
\renewcommand{\headrulewidth}{0.4pt}
% A 12pt, fancyhdr reclame 13,6pt de bande d'en-tete et avertit a chaque page.
\setlength{\headheight}{14pt}

% --- RACCOURCIS ---
\newcommand{\code}[1]{\texttt{\small #1}}
\newcommand{\hubeaudepot}{\code{hubeau\_data\_integration}}
\newcommand{\plateforme}{\code{time-serie-explo}}
